\documentclass[aspectratio=169,spanish,xcolor=table]{beamer}
\usepackage{iftex}
\ifPDFTeX
  \usepackage[utf8]{inputenc}
  \usepackage[T1]{fontenc}
\fi
\usepackage[spanish,es-tabla,es-noshorthands]{babel}
\usepackage{helvet}
\renewcommand{\familydefault}{\sfdefault}
\usepackage{tikz,pgfplots,array,booktabs,ragged2e}
\pgfplotsset{compat=1.18}
\usetikzlibrary{arrows.meta,positioning,calc,shapes.geometric,fit}

\definecolor{iiiepeNavy}{HTML}{163B5C}
\definecolor{iiiepeBlue}{HTML}{1F6F9F}
\definecolor{iiiepeAqua}{HTML}{20B7A8}
\definecolor{iiiepeOrange}{HTML}{F28C28}
\definecolor{iiiepePaper}{HTML}{F6F8F7}
\definecolor{iiiepeInk}{HTML}{1E2A32}
\definecolor{iiiepePale}{HTML}{EEF5F8}
\definecolor{iiiepeMint}{HTML}{E7F6F4}
\definecolor{iiiepeWarm}{HTML}{FFF3E6}

\usetheme{default}
\usecolortheme{default}
\setbeamertemplate{navigation symbols}{}
\setbeamercolor{normal text}{fg=iiiepeInk,bg=white}
\setbeamercolor{structure}{fg=iiiepeNavy}
\setbeamercolor{frametitle}{fg=white,bg=iiiepeNavy}
\setbeamercolor{block title}{fg=white,bg=iiiepeBlue}
\setbeamercolor{block body}{fg=iiiepeInk,bg=iiiepePale}
\setbeamercolor{item}{fg=iiiepeOrange}

\setbeamertemplate{frametitle}{%
  \nointerlineskip%
  \begin{beamercolorbox}[wd=\paperwidth,ht=0.95cm,dp=0.22cm,leftskip=0.45cm,rightskip=0.45cm]{frametitle}%
    \begin{minipage}[c]{13.2cm}\usebeamerfont{frametitle}\large\bfseries\insertframetitle\end{minipage}%
    \hfill\begin{minipage}[c]{1.6cm}\raggedleft\includegraphics[height=0.56cm]{img/iiiepe/logo-iiiepe-blanco.png}\end{minipage}%
  \end{beamercolorbox}%
}

\setbeamertemplate{footline}{%
  \leavevmode%
  \hbox{%
    \begin{beamercolorbox}[wd=.78\paperwidth,ht=0.5cm,dp=0.18cm,leftskip=0.45cm]{frametitle}%
      \scriptsize\color{white}M. de la Cruz Muñoz · MAEM · Matrícula 00874 · SD26 · Actividad 1 · Sesión 1%
    \end{beamercolorbox}%
    \begin{beamercolorbox}[wd=.22\paperwidth,ht=0.5cm,dp=0.18cm,rightskip=0.45cm]{frametitle}%
      \hfill\scriptsize\color{white}\insertframenumber{} / \inserttotalframenumber%
    \end{beamercolorbox}%
  }%
  \vskip0pt%
}

\newcommand{\Pill}[2]{%
  \tikz[baseline=(x.base)]{\node[rounded corners=3pt,fill=#1!14,draw=#1!75!black,thick,inner xsep=6pt,inner ysep=4pt,align=center] (x) {\scriptsize\bfseries #2};}%
}

\title{El Futuro del Aprendizaje y Diagnóstico Matemático}
\subtitle{Actividad 1 · Sesión 1 · Evaluación Diagnóstica y Análisis de Cynthia Luna Scott (2015)}
\author{Mtro. Martín Jonathan de la Cruz Muñoz · Matrícula 00874}
\institute{IIIEPE · Maestría en Aprendizaje y Enseñanza de las Matemáticas · Grupo A}
\date{19 de noviembre de 2026}

\begin{document}

\begin{frame}[plain]
  \begin{tikzpicture}[remember picture,overlay]
    \fill[iiiepeNavy] (current page.north west) rectangle ([xshift=10.3cm]current page.south west);
    \fill[iiiepeBlue] ([xshift=10.3cm]current page.north west) rectangle (current page.south east);
    \fill[iiiepeOrange] ([xshift=0.8cm,yshift=-6.45cm]current page.north west) rectangle ([xshift=9.3cm,yshift=-6.53cm]current page.north west);
    \node[anchor=north west,text width=8.7cm,align=left] at ([xshift=0.8cm,yshift=-1.15cm]current page.north west) {%
      {\color{iiiepeAqua!45!white}\small\bfseries MAESTRÍA EN APRENDIZAJE Y ENSEÑANZA DE LAS MATEMÁTICAS · SD26\par}%
      \vspace{0.35cm}
      {\color{white}\Large\bfseries \inserttitle\par}%
      \vspace{0.25cm}
      {\color{white!90}\normalsize \insertsubtitle\par}%
    };
    \node[anchor=south west,text width=8.7cm,align=left] at ([xshift=0.8cm,yshift=0.75cm]current page.south west) {%
      {\color{white}\normalsize\bfseries Mtro. Martín Jonathan de la Cruz Muñoz\par}%
      {\color{white!85}\scriptsize Matrícula 00874 · martin.delacruz@campus.iiiepe.edu.mx\par}%
      {\color{iiiepeAqua!35!white}\scriptsize Monterrey, Nuevo León · 19 de noviembre de 2026\par}%
    };
    \node[anchor=center,text width=5.0cm,align=center] at ([xshift=13.15cm,yshift=-4.5cm]current page.north west) {%
      \includegraphics[width=3.6cm]{img/iiiepe/logo-iiiepe-blanco.png}\par\vspace{0.45cm}%
      {\color{white}\large\bfseries IIIEPE\par}%
      \vspace{0.1cm}%
      {\color{white!85}\scriptsize Fundamentos para la enseñanza\\y el aprendizaje I\par}%
    };
  \end{tikzpicture}
\end{frame}
\begin{frame}{Propósito y Tesis Central (Sesión 1)}
  \begin{block}{Planteamiento de partida}
    Enseñar matemáticas en bachillerato técnico e ingeniería exige superar la mecanización sin contexto para construir comprensión conceptual, pensamiento crítico y sentido humano \textbf{(Scott, 2015; UNESCO, 2015)}.
  \end{block}
  \vspace{0.3cm}
  \begin{columns}[T]
    \column{0.48\textwidth}
      \textbf{\color{iiiepeNavy}Fuerzas de cambio (UNESCO):}
      \begin{itemize}\scriptsize
        \item Globalización y mutación del mercado laboral.
        \item Transformación tecnológica acelerada.
        \item Cambios demográficos y movilidad profesional.
        \item Brechas económicas y urgencia ambiental.
      \end{itemize}
    \column{0.48\textwidth}
      \textbf{\color{iiiepeBlue}Ejes Conceptuales de la Sesión 1:}
      \begin{itemize}\scriptsize
        \item Selección y análisis propio de 4 ideas centrales de Cynthia Luna Scott (2015).
        \item Diagnóstico de cuatro conceptos clave: enseñanza, aprendizaje, paradigmas y matemáticas.
        \item Ilustración matemática situada para docencia técnica.
      \end{itemize}
  \end{columns}
\end{frame}

\begin{frame}{Cuatro Ideas Relevantes de Cynthia Luna Scott (2015) · I}
  \begin{columns}[T]
    \column{0.48\textwidth}
      \begin{block}{Idea 1 · Crisis del modelo transmisivo}
        La memorización para exámenes estandarizados es insuficiente ante problemas reales complejos.\\[0.2cm]
        \textbf{Comentario:} La agilidad al despejar $x$ puede ocultar incapacidad para identificar variables, tolerancias o supuestos físicos en taller.
      \end{block}
    \column{0.48\textwidth}
      \begin{block}{Idea 2 · El nuevo aprendiz digital}
        El acceso ubicuo a pantallas y redes no garantiza por sí solo aprendizaje crítico ni autorregulación.\\[0.2cm]
        \textbf{Comentario:} Usar graficadores o software sin comprender tasas de cambio ni unidades acelera el error si falta mediación docente.
      \end{block}
  \end{columns}
\end{frame}

\begin{frame}{Cuatro Ideas Relevantes de Cynthia Luna Scott (2015) · II}
  \begin{columns}[T]
    \column{0.48\textwidth}
      \begin{block}{Idea 3 · Brecha escuela--mundo laboral}
        Las tareas rutinarias disminuyen; se demandan pensamiento experto y solución de problemas no estructurados.\\[0.2cm]
        \textbf{Comentario:} En mantenimiento y mecatrónica los problemas llegan como síntomas reales; la matemática escolar debe habilitar el modelado bajo incertidumbre.
      \end{block}
    \column{0.48\textwidth}
      \begin{block}{Idea 4 · Pedagogía 2.0 y ecosistemas}
        Convergencia de los cuatro pilares de Delors (1996) en redes de aprendizaje formales e informales.\\[0.2cm]
        \textbf{Comentario:} Digitalizar sin dialogar no cambia la escuela; conectar aula, taller y comunidad fortalece equidad y sentido.
      \end{block}
  \end{columns}
\end{frame}

\begin{frame}{Nexo Diagnóstico: Cuatro Conceptos Fundamentales}
  \centering
  \begin{tikzpicture}[
    every node/.style={thick,rounded corners=4pt,text width=5.4cm,align=left,inner sep=6pt,font=\scriptsize},
    arr/.style={->,thick,>=Latex,iiiepeNavy}
  ]
    \node[draw=iiiepeNavy,fill=iiiepeNavy!12] (a) at (0,2.4) {\textbf{Enseñanza:} Mediación intencional que genera preguntas, andamiajes y condiciones de participación equitativa.};
    \node[draw=iiiepeBlue,fill=iiiepeBlue!12] (b) at (7.0,2.4) {\textbf{Aprendizaje:} Construcción activa y relacional que integra comprensión, diálogo y transferencia.};
    \node[draw=iiiepeAqua!75!black,fill=iiiepeAqua!12] (c) at (0,0) {\textbf{Paradigmas:} Marcos epistemológicos sobre qué vale como conocimiento, cómo se trata el error y para qué se evalúa.};
    \node[draw=iiiepeOrange!85!black,fill=iiiepeOrange!12] (d) at (7.0,0) {\textbf{Matemáticas:} Lenguaje de variación, estructura e inferencia para decidir con evidencia.};
    \draw[arr] (a) -- (b);
    \draw[arr] (c) -- (a);
    \draw[arr] (c) -- (d);
    \draw[arr] (d) -- (b);
  \end{tikzpicture}
\end{frame}

\begin{frame}{Del Paradigma Transmisivo al Enfoque Comprensivo}
  \scriptsize\renewcommand{\arraystretch}{1.4}
  \begin{tabular}{p{2.3cm}p{5.4cm}p{5.4cm}}
    \toprule
    \rowcolor{iiiepePale} \textbf{Eje} & \textbf{Práctica transmisiva tradicional} & \textbf{Horizonte formativo (SD26)} \\
    \midrule
    Enseñanza & Exposición de reglas aisladas y copia de ejercicios tipo. & Diseño de retos situados, indagación guiada y diálogo entre pares. \\
    Aprendizaje & Memorización efímera para acreditar evaluaciones sumativas. & Apropiación conceptual capaz de explicar el porqué del procedimiento. \\
    Error y evaluación & Penalización inmediata; enfoque puramente calificativo. & Evidencia diagnóstica para ajustar el andamiaje en el grupo. \\
    Matemáticas & Manipulación simbólica desligada de unidades y realidad. & Modelación crítica aplicada a sistemas técnicos y vida comunitaria. \\
    \bottomrule
  \end{tabular}
\end{frame}

\begin{frame}{Ejemplos Matemáticos: Lectura Relacional vs Mecanización}
  \begin{columns}[T]
    \column{0.48\textwidth}
      \begin{block}{1. Cinemática lineal en taller}
        $$d(t) = 0.65\,t + 1.4$$
        \scriptsize
        \begin{itemize}
          \item \textbf{Cálculo:} $d(8) = 6.6\text{ m}$.
          \item \textbf{Comprensión profunda:} Distinguir función afín vs proporcional ($d_0 = 1.4\text{ m}$), interpretar la pendiente ($0.65\text{ m/s}$) y verificar límites del riel físico.
        \end{itemize}
      \end{block}
    \column{0.48\textwidth}
      \begin{block}{2. Costos de servicio técnico}
        $$C_A(h) = 250 + 80h, \quad C_B(h) = 100 + 110h$$
        \scriptsize
        \begin{itemize}
          \item \textbf{Equilibrio:} $h = 5\text{ horas}$ (\$650).
          \item \textbf{Decisión argumentada:} Por debajo de $5\text{ h}$ conviene $B$; para jornadas mayores a $5\text{ h}$ conviene $A$ por su menor costo marginal (\$80/h).
        \end{itemize}
      \end{block}
  \end{columns}
\end{frame}

\begin{frame}{Delors (1996) y Ecosistemas en la Enseñanza Técnica}
  \centering
  \Pill{iiiepeNavy}{Aprender a conocer} \hspace{0.2cm}
  \Pill{iiiepeBlue}{Aprender a hacer} \hspace{0.2cm}
  \Pill{iiiepeAqua!75!black}{Aprender a convivir} \hspace{0.2cm}
  \Pill{iiiepeOrange!85!black}{Aprender a ser}
  \vspace{0.55cm}
  \begin{block}{Puente con la docencia en CONALEP, CECyTE e Ingeniería}
    \begin{itemize}\small
      \item Vincular el aula de matemáticas con el taller de mantenimiento, el análisis de sensores y problemas de la comunidad.
      \item Exigir justificación de unidades, validez del modelo y lectura crítica de gráficas generadas digitalmente.
      \item Cuidar la equidad mediante apoyos graduados para estudiantes con trayectorias previas desiguales.
    \end{itemize}
  \end{block}
\end{frame}

\begin{frame}{Ruta de Transformación de la Práctica Docente}
  \centering
  \begin{tikzpicture}[
    node distance=0.8cm,
    s/.style={draw=iiiepeNavy,fill=iiiepePale,rounded corners=4pt,text width=2.85cm,minimum height=1.7cm,align=center,font=\scriptsize},
    arr/.style={->,very thick,iiiepeOrange,>=Latex}
  ]
    \node[s] (s1) {\textbf{1. Diagnóstico}\\Reconocer ideas previas y bloqueos reales.};
    \node[s,right=of s1] (s2) {\textbf{2. Problematizar}\\Plantear retos técnicos o comunitarios auténticos.};
    \node[s,right=of s2] (s3) {\textbf{3. Modelar}\\Conectar tablas, gráficas y álgebra con sentido.};
    \node[s,right=of s3] (s4) {\textbf{4. Dialogar}\\Argumentar errores, límites y decisiones.};
    \draw[arr] (s1) -- (s2);
    \draw[arr] (s2) -- (s3);
    \draw[arr] (s3) -- (s4);
  \end{tikzpicture}
\end{frame}

\begin{frame}{Conclusiones de la Sesión 1}
  \begin{itemize}\normalsize
    \item Las seis fuerzas de transformación obligan a abandonar la enseñanza centrada únicamente en memorizar algoritmos \textbf{(Scott, 2015)}.
    \item Enseñar matemáticas en contextos técnicos requiere conjugar exactitud instrumental, modelación crítica y diálogo inclusivo.
    \item La Tecnología solo potencia el aprendizaje cuando está gobernada por preguntas conceptuales e interpretación de evidencia.
    \item Este diagnóstico establece la base reflexiva para examinar la propia identidad docente en la Sesión 2.
  \end{itemize}
\end{frame}

\begin{frame}{Referencias Principales (APA 7)}
  \scriptsize\setlength{\parskip}{6pt}
  \begin{itemize}
    \item CONALEP. (2025). \textit{Perfil de egreso: Profesional Técnico-Bachiller en Mecatrónica}. Secretaría Académica.
    \item Delors, J., et al. (1996). \textit{La educación encierra un tesoro}. Santillana/UNESCO.
    \item Pimienta Prieto, J. H. (2012). \textit{Estrategias de enseñanza-aprendizaje}. Pearson Educación.
    \item Sawyer, R. K. (2008). Optimizing learning: Implications of learning sciences research. En \textit{Innovating to learn, learning to innovate} (pp. 45--65). OECD.
    \item Scott, C. L. (2015). \textit{El futuro del aprendizaje 1}. UNESCO.
    \item SEP. (2023). \textit{Programa de estudios: Pensamiento Matemático I}. SEMS.
    \item UNESCO. (2015). \textit{Replantear la educación: ¿Hacia un bien común mundial?} UNESCO.
  \end{itemize}
\end{frame}

\end{document}
